\chapter{Carbonyls: Nucleophilic Additions, Acetals and Protection}\label{ch:b1:acetals-protection}

Dissolve glucose in water and almost none of it stays as the open-chain
aldehyde drawn in textbooks: the molecule bends round, its own hydroxyl
group adds to its aldehyde, and a six-membered ring closes. The same
addition of an alcohol to a carbonyl group, repeated once more with an
acid \omterm{def:b1:elementary-steps:catalyst}{catalyst}, gives \omterm{def:b1:acetals-protection:hemiacetal}{acetals}, compounds stable towards bases and the most
reactive \omterm{def:b1:electronic-effects:nucleophile}{nucleophiles}. That stability is a tool: an aldehyde that would
be destroyed by a \omterm{def:b1:organomagnesium:organometallic}{Grignard reagent} can be hidden as an \omterm{def:b1:acetals-protection:hemiacetal}{acetal}, carried
through the reaction, and recovered at the end. This chapter studies the
carbonyl group as an \omterm{def:b1:electronic-effects:nucleophile}{electrophile}, the additions of water and alcohols,
and the strategy of \omterm{def:b1:acetals-protection:protecting-group}{protection}.

\begin{recall}
Grignard additions to carbonyl compounds (\cref{ch:b1:organomagnesium});
the \omterm{def:b1:extent-q-and-k:quotient}{reaction quotient} $Q$ and the evolution of a system while $Q \neq K$
(\cref{ch:b1:extent-q-and-k}); \omterm{def:b1:electronic-effects:curly-arrow}{curly arrows}, \omterm{def:b1:electronic-effects:nucleophile}{nucleophiles} and
\omterm{def:b1:electronic-effects:nucleophile}{electrophiles} (\cref{ch:b1:electronic-effects}). The school volume (grade
12) introduced the idea of protecting a group during a synthesis.
\end{recall}

\section{The carbonyl group as an electrophile}

The C=O bond is polarised towards oxygen (\cref{ch:b1:periodicity}), and
a \omterm{def:b1:lewis-resonance-vsepr:resonance}{resonance structure} \ce{C+-O-} puts a full positive charge on carbon:
the carbon is electrophilic, attacked by \omterm{def:b1:electronic-effects:nucleophile}{nucleophiles} perpendicular to the
plane of the group. Strong \omterm{def:b1:electronic-effects:nucleophile}{nucleophiles} (\omterm{def:b1:organomagnesium:organometallic}{Grignard reagents}, hydride,
cyanide) add directly. Weak \omterm{def:b1:electronic-effects:nucleophile}{nucleophiles} --- water, alcohols --- need help.

\begin{definition}[Electrophilic activation]\label{def:b1:acetals-protection:activation}
\emph{Electrophilic activation}\index{electrophilic activation} of a
carbonyl group is its protonation (or its coordination to a \omterm{def:b1:lewis-resonance-vsepr:lewis-acid}{Lewis acid}) on
oxygen. The cation \ce{C=OH+}, whose positive charge is shared by carbon
through the \omterm{def:b1:lewis-resonance-vsepr:resonance}{resonance structure} \ce{C+-OH}, is attacked even by weak
\omterm{def:b1:electronic-effects:nucleophile}{nucleophiles}.
\end{definition}

\begin{omfigure}
\setchemfig{atom sep=2.1em}
\schemestart
\chemfig{C(-[:150]R)(-[:210]H)=O}
\arrow{->[\ce{H+}]}
\chemfig{C(-[:150]R)(-[:210]H)=O^{+}-H}
\arrow{<->}
\chemfig{C^{+}(-[:150]R)(-[:210]H)-O-H}
\schemestop

{\small Activation of an aldehyde by acid: protonation of the oxygen
produces a cation whose \omterm{def:b1:lewis-resonance-vsepr:resonance}{resonance structure} carries the charge on
carbon.}
\end{omfigure}

\section{Hydrates and hemiacetals}

Water adds reversibly to aldehydes and ketones, \ce{R2C=O + H2O <=>
R2C(OH)2}; the hydrate is usually minor for ketones and can be major for
methanal. An alcohol adds the same way.

\begin{definition}[Hemiacetal and acetal]\label{def:b1:acetals-protection:hemiacetal}
A \emph{hemiacetal}\index{hemiacetal} is a compound in which one carbon
carries both an OH and an OR group, formed by the addition of an alcohol
to a carbonyl group. An \emph{acetal}\index{acetal} is a compound in which
one carbon carries two OR groups, formed from a hemiacetal and a second
alcohol molecule, with loss of water.
\end{definition}

Open-chain \omterm{def:b1:acetals-protection:hemiacetal}{hemiacetals} are usually minor in their equilibrium with the
aldehyde and the alcohol; when the OH and the C=O belong to the same
molecule and a five- or six-membered ring can close, the cyclic
\omterm{def:b1:acetals-protection:hemiacetal}{hemiacetal} predominates. That is the case of glucose, and of the simpler
5-hydroxypentanal.

\begin{omfigure}
\setchemfig{atom sep=2em}
\schemestart
\chemfig{HO-[:30]-[:-30]-[:30]-[:-30]-[:30]C(=[:90]O)-[:-30]H}
\arrow{<=>}
\chemfig{*6(O-(-[:-30]OH)-----)}
\schemestop

\medskip

{\small 5-Hydroxypentanal and its cyclic \omterm{def:b1:acetals-protection:hemiacetal}{hemiacetal}, a six-membered ring
containing the oxygen. The OH of carbon 5 has added to the aldehyde of
carbon 1; the ring form predominates, as for glucose.}
\end{omfigure}

\section{Acetals}

\begin{proposition}[Mechanism of acetal formation]\label{prop:b1:acetals-protection:acetal-mechanism}
In acid, an aldehyde or ketone and an alcohol give the \omterm{def:b1:acetals-protection:hemiacetal}{hemiacetal}, then
the \omterm{def:b1:acetals-protection:hemiacetal}{acetal}, by the steps: protonation of C=O; addition of the alcohol;
loss of a proton (\omterm{def:b1:acetals-protection:hemiacetal}{hemiacetal}); protonation of the OH; loss of water,
giving a \omterm{def:b1:electronic-effects:intermediates}{carbocation} stabilised by the remaining OR group (an oxocarbenium
ion); addition of a second alcohol; loss of a proton. Every step is
reversible.
\end{proposition}

\begin{proof}
Each step is an acid--base proton transfer or the addition (or loss) of a
\omterm{def:b1:electronic-effects:nucleophile}{nucleophile} to (or from) an electrophilic carbon, as described in
\cref{ch:b1:electronic-effects}. The loss of water is possible because
the cation formed, $\mathrm{R_2C^+{-}OR}$, has a \omterm{def:b1:lewis-resonance-vsepr:resonance}{resonance structure} $\mathrm{R_2C{=}O^+R}$
in which every atom has an octet (\cref{prop:b1:electronic-effects:carbocation-order}).
Acid is consumed in the protonations and regenerated in the
deprotonations: it is a \omterm{def:b1:elementary-steps:catalyst}{catalyst}.
\end{proof}

\begin{omfigure}
\setchemfig{atom sep=1.9em}
\schemestart
\chemfig{C(-[:150]R)(-[:210]H)=O}
\arrow{->[\ce{H+}][]}
\chemfig{C^{+}(-[:150]R)(-[:210]H)-OH}
\arrow{->[$\mathrm{R'OH}$][\ce{-H+}]}
\chemfig{C(-[:150]R)(-[:210]H)(-[2]OR')-OH}
\schemestop

\medskip
\schemestart
\chemfig{C(-[:150]R)(-[:210]H)(-[2]OR')-OH}
\arrow{->[\ce{H+}][\ce{-H2O}]}
\chemfig{C^{+}(-[:150]R)(-[:210]H)-OR'}
\arrow{->[$\mathrm{R'OH}$][\ce{-H+}]}
\chemfig{C(-[:150]R)(-[:210]H)(-[2]OR')-OR'}
\schemestop
\par\vspace{6pt}

{\small \omterm{def:b1:acetals-protection:hemiacetal}{Acetal} formation, in two lines: first the \omterm{def:b1:acetals-protection:hemiacetal}{hemiacetal} (activation,
addition of the alcohol, loss of a proton), then the \omterm{def:b1:acetals-protection:hemiacetal}{acetal} (protonation of
the OH, loss of water to the stabilised cation, addition of a second
alcohol, loss of a proton).}
\end{omfigure}

The overall reaction, $\mathrm{RCHO} + 2\,\mathrm{R'OH} \rightleftharpoons \mathrm{RCH(OR')_2} + \ce{H2O}$, has a
modest equilibrium constant, and water is a product.

\begin{proposition}[Driving the acetalisation]\label{prop:b1:acetals-protection:dean-stark}
If the water formed is removed continuously, the acetalisation goes on
until the carbonyl compound (or the alcohol) is used up.
\end{proposition}

\begin{proof}
The \omterm{def:b1:extent-q-and-k:quotient}{reaction quotient} $Q = [\text{acetal}][\ce{H2O}]/([\text{aldehyde}]
[\mathrm{R'OH}]^2)$ stays below $K$ as long as water is removed: by
\cref{ch:b1:extent-q-and-k} the reaction keeps going forward, whatever the
value of $K$. With a diol (ethane-1,2-diol) the \omterm{def:b1:acetals-protection:hemiacetal}{acetal} is cyclic, and the
equilibrium is also more favourable, one molecule of diol replacing two of
alcohol.
\end{proof}

\begin{omfigure}
\begin{tikzpicture}[font=\small]
  \pic at (0,0) {heatingmantle};
  \pic at (0,0) {rbflask={0.4}{omDef!14}};
  \pic (d) at (0,3.0) {deanstark={0.35}{omDef!35}{orange!15}};
  \pic at (0,6.4) {condenser};
  \draw[->] (2.4,4.0) node[right, align=left] {collector: water (lower layer)\\returns no further; toluene\\overflows back to the flask} -- (0.8,3.8);
  \draw[->] (-2.2,1.0) node[left, align=right] {aldehyde, diol, toluene,\\acid catalyst, boiling} -- (-0.7,0.9);
  \draw[->] (1.8,8.0) node[right] {condenser} -- (0.4,7.8);
\end{tikzpicture}

{\small A Dean--Stark trap. The vapours of toluene and water condense and
fall into the graduated collector; water, denser and not \omterm{def:b1:intermolecular-forces-solvents:miscibility}{miscible} with
toluene, sinks and stays, toluene overflows back into the flask. The
volume of water collected measures the progress of the reaction.}
\end{omfigure}

\begin{proposition}[Stability of acetals]\label{prop:b1:acetals-protection:acetal-stability}
\omterm{def:b1:acetals-protection:hemiacetal}{Acetals} are stable in neutral and basic media and towards \omterm{def:b1:electronic-effects:nucleophile}{nucleophiles},
\omterm{def:b1:organomagnesium:organometallic}{Grignard reagents} and hydride reagents; aqueous acid hydrolyses them back
to the carbonyl compound and the alcohol.
\end{proposition}

\begin{proof}
An \omterm{def:b1:acetals-protection:hemiacetal}{acetal} carbon carries no \omterm{def:b1:electronic-effects:nucleophile}{leaving group} that a base or a \omterm{def:b1:electronic-effects:nucleophile}{nucleophile}
could expel: \ce{RO-}, a \omterm{def:b1:predominant-reaction:strong-acid}{strong base}, does not leave, and there is no C=O
left to attack. In aqueous acid the mechanism above runs backwards: the
large excess of water makes $Q < K$ for the reverse reaction.
\end{proof}

\section{Protecting and deprotecting}

\begin{definition}[Protecting group]\label{def:b1:acetals-protection:protecting-group}
A \emph{protecting group}\index{protecting group} is a group introduced
into a molecule to mask a function temporarily, so that it does not react
in a step aimed at another part of the molecule. Its introduction is the
\emph{protection}\index{protection}, its removal the
\emph{deprotection}\index{deprotection}.
\end{definition}

\begin{method}[Protect, react, deprotect]\label{met:b1:acetals-protection:protect}
\begin{enumerate}
  \item Identify the function that would react, or would destroy the
        reagent, in the planned step.
  \item Choose a \omterm{def:b1:acetals-protection:protecting-group}{protecting group} stable in that step and removable in
        conditions the rest of the molecule tolerates (for an aldehyde or
        ketone facing bases, \omterm{def:b1:electronic-effects:nucleophile}{nucleophiles} or hydrides: a cyclic \omterm{def:b1:acetals-protection:hemiacetal}{acetal}).
  \item Protect; carry out the planned step; deprotect.
  \item Count the cost: two more steps, each with its \omterm{def:b1:extent-q-and-k:fractional-extent}{yield}.
\end{enumerate}
\end{method}

\begin{example}[A ketone in the way]\label{ex:b1:acetals-protection:ketoester}
To add a \omterm{def:b1:organomagnesium:organometallic}{Grignard reagent} to the aldehyde group of a molecule that also
carries a ketone, no simple \omterm{def:b1:acetals-protection:hemiacetal}{acetal} choice exists (aldehydes form \omterm{def:b1:acetals-protection:hemiacetal}{acetals}
more readily than ketones, so the wrong group would be protected); the
order of the steps or the reagents must change. \omterm{def:b1:acetals-protection:protecting-group}{Protection} is a
strategy, not a reflex.
\end{example}

\section{Exercises}

\begin{exercise}[$\star$]\label{exo:b1:acetals-protection:1}
Identify \omterm{def:b1:acetals-protection:hemiacetal}{hemiacetal}, \omterm{def:b1:acetals-protection:hemiacetal}{acetal}, hydrate, ether or alcohol carbons in:
\ce{CH3CH(OH)OCH3}, \ce{CH3CH(OCH3)2}, \ce{CH3CH(OH)2}, \ce{CH3OCH3},
\ce{(CH3)2C(OCH2CH3)2}.
\end{exercise}

\begin{exercise}[$\star$]\label{exo:b1:acetals-protection:2}
Write the overall equation of the formation of the \omterm{def:b1:acetals-protection:hemiacetal}{acetal} of propanal with
methanol, and of the cyclic \omterm{def:b1:acetals-protection:hemiacetal}{acetal} of propanone with ethane-1,2-diol.
\end{exercise}

\begin{exercise}[$\star$]\label{exo:b1:acetals-protection:3}
Which of these reagents leave an \omterm{def:b1:acetals-protection:hemiacetal}{acetal} unchanged: sodium hydroxide
solution, methylmagnesium bromide, dilute sulfuric acid in water, sodium
borohydride?
\end{exercise}

\begin{exercise}[$\star$]\label{exo:b1:acetals-protection:4}
Draw the cyclic \omterm{def:b1:acetals-protection:hemiacetal}{hemiacetal} formed by 4-hydroxybutanal. What is the size
of the ring?
\end{exercise}

\begin{exercise}[$\star\star$]\label{exo:b1:acetals-protection:5}
Write the full mechanism of the acid-catalysed formation of the dimethyl
\omterm{def:b1:acetals-protection:hemiacetal}{acetal} of ethanal, with \omterm{def:b1:electronic-effects:curly-arrow}{curly arrows}, and show that the acid is
regenerated.
\end{exercise}

\begin{exercise}[$\star\star$]\label{exo:b1:acetals-protection:6}
Write the mechanism of the acid hydrolysis of
2,2-dimethoxypropane. Why is a large amount of water used?
\end{exercise}

\begin{exercise}[$\star\star$]\label{exo:b1:acetals-protection:7}
An acetalisation of \qty{0.250}{mol} of a ketone is followed with a
Dean--Stark trap. What volume of water should be collected at full
conversion (density about \qty{1.0}{g/mL})? After an hour,
\qty{3.2}{mL} have collected: what is the conversion?
\end{exercise}

\begin{exercise}[$\star\star$]\label{exo:b1:acetals-protection:8}
Explain why a cyclic \omterm{def:b1:acetals-protection:hemiacetal}{acetal} with ethane-1,2-diol forms more completely
than the \omterm{def:b1:acetals-protection:hemiacetal}{acetal} with two molecules of methanol, for the same carbonyl
compound.
\end{exercise}

\begin{exercise}[$\star\star$]\label{exo:b1:acetals-protection:9}
Why must the acid \omterm{def:b1:elementary-steps:catalyst}{catalyst} be removed (or neutralised) before a protected
compound is treated with a \omterm{def:b1:organomagnesium:organometallic}{Grignard reagent}?
\end{exercise}

\begin{exercise}[$\star\star\star$]\label{exo:b1:acetals-protection:10}
Propose a synthesis of 5-hydroxy-5-methylhexan-2-one,
\ce{(CH3)2C(OH)CH2CH2COCH3}, from 4-chlorobutan-2-one and propanone, using
a \omterm{def:b1:acetals-protection:protecting-group}{protection}. Give each step, its reagents and the role of each.
\end{exercise}

\begin{exercise}[$\star\star\star$]\label{exo:b1:acetals-protection:11}
Show with the \omterm{def:b1:extent-q-and-k:quotient}{reaction quotient} that a large excess of methanol also
drives the acetalisation of an aldehyde, and compare this method with the
removal of water.
\end{exercise}

\begin{exercise}[$\star\star\star$]\label{exo:b1:acetals-protection:12}
Explain why the cyclic \omterm{def:b1:acetals-protection:hemiacetal}{hemiacetal} of glucose is attacked, in acidic
methanol, to give a methyl \omterm{def:b1:acetals-protection:hemiacetal}{acetal} (a methyl glycoside) while the other OH
groups of glucose are not converted into ethers.
\end{exercise}

\section{Problem: A Grignard Reagent that Carries an Aldehyde}

\begin{problem}[{Weekend problem --- why 4-bromobutanal cannot give a
Grignard reagent, its protection as a cyclic acetal with a Dean--Stark
trap, the Grignard addition to propanone, the deprotection, and the
volume of water collected}]
\label{pb:b1:acetals-protection:1}
A chemist needs 5-hydroxy-5-methylhexanal, \ce{(CH3)2C(OH)CH2CH2CH2CHO}, and
plans to make it from 4-bromobutanal, \ce{BrCH2CH2CH2CHO}, and propanone.
She starts from \qty{15.09}{g} of 4-bromobutanal. Molar masses
(\unit{g/mol}): \ce{C4H7BrO} 150.9, \ce{C2H6O2} 62.0, \ce{C6H11BrO2} 194.9,
\ce{C3H6O} 58.0, \ce{C9H18O3} 174.0, \ce{C7H14O2} 130.0, \ce{H2O} 18.0;
density of water \qty{0.997}{g/mL}.
% ledger: rho:water

\medskip
\noindent\textbf{Part I --- The problem.}
\begin{enumerate}
  \item Write the \omterm{def:b1:organomagnesium:organometallic}{Grignard reagent} that 4-bromobutanal would give.
  \item Why can that reagent not exist? Write the reaction that destroys
        it.
  \item Which function must be protected, and against what?
  \item Why is an \omterm{def:b1:acetals-protection:hemiacetal}{acetal} suitable?
  \item Why is a cyclic \omterm{def:b1:acetals-protection:hemiacetal}{acetal} (with ethane-1,2-diol) preferred?
  \item Compute the amount of 4-bromobutanal.
  \item Compute the mass of diol for a \qty{20}{\percent} excess.
\end{enumerate}

\medskip
\noindent\textbf{Part II --- \omterm{def:b1:acetals-protection:protecting-group}{Protection}.}
\begin{enumerate}[resume]
  \item Write the equation of the acetalisation with ethane-1,2-diol.
  \item What is the role of the acid \omterm{def:b1:elementary-steps:catalyst}{catalyst} (4-methylbenzenesulfonic
        acid)?
  \item Describe the operation of the Dean--Stark trap.
  \item Explain with $Q$ and $K$ why removing water drives the reaction.
  \item Why is the toluene, and not the water, returned to the flask?
  \item How does the chemist know the reaction is complete?
  \item Compute the mass of protected bromide expected at full conversion.
\end{enumerate}

\medskip
\noindent\textbf{Part III --- The Grignard step.}
\begin{enumerate}[resume]
  \item Write the formation of the \omterm{def:b1:organomagnesium:organometallic}{Grignard reagent} from the protected
        bromide.
  \item Write its addition to propanone and the \omterm{def:b1:alcohol-activation:alkoxide}{alkoxide} formed.
  \item Why must the work-up of this step be mildly basic or neutral (for
        example aqueous ammonium chloride) rather than acidic?
  \item Compute the mass of protected alcohol for a \qty{70}{\percent} \omterm{def:b1:extent-q-and-k:fractional-extent}{yield}
        of this step.
\end{enumerate}

\medskip
\noindent\textbf{Part IV --- \omterm{def:b1:acetals-protection:protecting-group}{Deprotection}.}
\begin{enumerate}[resume]
  \item Write the \omterm{def:b1:acetals-protection:protecting-group}{deprotection} in aqueous acid.
  \item Why does the tertiary alcohol survive mild aqueous acid?
  \item The product can close a six-membered ring by adding its OH to the
        aldehyde. Draw that cyclic \omterm{def:b1:acetals-protection:hemiacetal}{hemiacetal}.
  \item Compute the mass of hydroxy aldehyde at \qty{90}{\percent} \omterm{def:b1:extent-q-and-k:fractional-extent}{yield} of
        \omterm{def:b1:acetals-protection:protecting-group}{deprotection}.
  \item What overall \omterm{def:b1:extent-q-and-k:fractional-extent}{yield} from 4-bromobutanal does the chemist reach?
  \item Compute the mass of water formed in the \omterm{def:b1:acetals-protection:protecting-group}{protection} step at full
        conversion.
  \item State the volume of water collected in the Dean--Stark trap at full
        conversion.
\end{enumerate}
\end{problem}
